\documentclass[12pt]{article}
\usepackage{seantex}

% --- Document Info ---
\title{ÜS Analysis II: Week 8, No. 2}
\author{Sean Findlay}
\date{\today}

% --- Start Document ---
\begin{document}

\maketitle

\section{Riemann-Integral in mehreren Variablen}
$Q = I_1 \times I_2 \times \dots, \ \lambda(Q) = \mu_n(Q) = |I_2| \cdot |I_2| \cdot \dots$. 
Treppenfunktion: $Q_1, \dots, Q_n$, $g_1, g_2, \dots, g_n \in \R$.

$$
g = \sum_{i = 1}^{n} g \cdot \mathds{1}_{Q_i}(x)
$$

Integral über Treppenfunktion:

$$
    \int_g = \sum g_i \cdot \lambda(Q_i) = \int_g \dd \lambda(x) = 2^{-np}\cdot\sum_i g_i
$$

\subsection*{Riemann-Integrierbar}

$f: A \subset \R^n \rightarrow \R$, 

$$
    I_{\text{low}} = \sup \left\{\int_g | g: \R^n \rightarrow \R \text{ Treppenfunktion}\right\}
    % I_{\text{high}} = \inf \left\{\int_g | g: \R^n \rightarrow \R \text{ Treppenfunktion}\right\}
$$

\begin{remark}{Eigenschaften des Riemann-Integrals}{}
    $f, g$ Treppenfunktionen, $\alpha \in \R$

    \begin{enumerate}
        \item Linearität:
        $$
            \int f + \alpha g = \int f + \alpha \cdot \int g 
        $$
        \item Dreiecksungleichung:
        $$
            |\int f| \leq \int |f|
        $$
        \item Monoton:
        $$
            f \leq g \implies \int f \leq \int g
        $$
    \end{enumerate}
\end{remark}

\begin{remark}{Integral über positive Funktionen}{}
    $f^+(x) = \max(f(x), 0), f^-(x) = \max(-f(x), 0)$,
    $$
        \int_A f = \int_A f^+ - \int_A f^-
    $$

    Wenn $f$ Riemann integrierbar, dann

    $$
        \int_A f(x) \dd\lambda(x) = \int_{a}^{b} f(x) \dd x
    $$
\end{remark}

\begin{remark}{Riemann-Integral und Jordan-Mass}{}
    $A \subset \R^n,\ f: A \rightarrow [0, \infty),\ \Gamma(f)_A = \{(x, y) \in \R^n \times \R \ |\ x \in A \land 0 \leq y \leq f(x)\}$ Hypograph von $f$. Dann gilt:
    
    $$
        \int f = \mu_{n + 1} \left(\Gamma_A(f) \right)
    $$
\end{remark}

\begin{theorem}{Fubini}{}
    $A_1 \subset \R^n, A_2 \subset \R^D$, $A_1 \times A_2$ beschränkt. $f: A \rightarrow [0, \infty)$ stetig.

    $$
        \int_{A_1 \times A_2} f \dd\lambda(x, y) = \int_{A_1} \left( \int_{A_2} f \dd \lambda(x)\right) \dd \lambda(y) = \int_{A_2} \left( \int_{A_1} f \dd \lambda(y)\right) \dd \lambda(x)
    $$
\end{theorem}

Wir integrieren unterschiedliche Variablen in beliebiger Reihenfolge.

\begin{example}{}{}
    $[0, 1] \times [0, 1],\ f(x, y) = 1$
    
    $$
        \int_{[0, 1]^2} f \dd \lambda(x, y) = \int_{0}^{1} x \int_{0}^{1} \dd y = 1
    $$
\end{example}

\begin{example}{}{}
    $A = \{(x, y) \in \R^2 \ | \ x \in [0, 1],\ x \leq y \leq e^x\}$ (Hypograph von $e^x$ im Intervall $[0, 1]$ mit $x \leq y$)

    $$
        V = \int_A x \cdot y \dd \lambda(x,y) = \int_A x \cdot y \mathds{1}_{[x, e^x]} \dd \lambda(x, y)
    $$

    wobei $A = [0, 1] \times [0, e]$

    $$
        = \int_{[0, 1]} \int_{[0, e]} x \cdot y \cdot \mathds{1}_{[x, e^x]} \cdot \dd y \dd x
    $$

    $$
        \int_{x}^{e^x} y \dd y = \left[\frac{y^2}{2}\right]_{x}^{e^x} = \frac{1}{2}e^x - \frac{x^2}{2}
    $$
\end{example}


\end{document}